% Step 5 appendix rework (2026-10-14): Appendix B = evaluation details (old App C, decoded). The protocol criteria
% tables, the controlled-environment validity rule, the claims table and the status table moved to the ESM (staging copy:
% esm_staging/old_app_b.tex, old_app_d.tex). Numbers unchanged; internal codes replaced by plain names.
\FloatBarrier
\section{Evaluation details}
\label{app:eval}

\subsection{Simulator and gate calibration}
\label{app:simparams}

The radio layer uses the path-loss models of 3GPP TR 38.901~\cite{3gpp38901}, with correlated shadowing and 3GPP
antenna patterns on a 20~MHz carrier. Interference is coupled to the previous tick's resource use, handover follows A3
with time-to-trigger, and a radio-link failure is detected with T310. A slice-aware scheduler serves a dedicated
low-latency pool, a work-conserving protected minimum share and a shared pool, and energy follows the EARTH power
model~\cite{auer2011earth}, with carrier shutdown for macro sectors and sleep for picos. The simulation tick is 40~ms,
key performance measurement (KPM) reports reach the RIC with a delay, and configuration changes take effect at the next
tick; each episode runs 120~s of warm-up and 600~s of scored time. The E2 interface is modelled only as KPM reports in
and control writes out, and no service-model conformance is claimed. The xApps act on delivered reports and the applied
configuration only, with per-seed variants of their dimensionless thresholds. The gate tolerances of
Equation~\eqref{eq:tol} used $G_{\mathrm{v}} = 19\,953.65$ violated UE-seconds and $G_{\mathrm{RLF}} = 35.025$ RLFs.

\subsection{Associational baselines}
\label{app:baselines}

Each baseline ran on the same unit table as PMRT and gave a score and a sign per hypothesis under the same support rule.
Correlation is the absolute Pearson correlation between the signed level and the change in the relation's KPI.
Granger~\cite{granger1969causal} regresses the post-window KPI on its pre-window value and the signed level and tests
the level's coefficient; it is reported with a tuned threshold and with BY control at $q = 0.05$ (Granger~+~BY).
SHAP-GBDT follows Sharma et al.~\cite{Sharma_2025}: it fits one gradient-boosted tree model per hypothesis and scores
the SHAP importance~\cite{lundberg2017shap} of the level, without the effect estimation of the original. The two-tower
network of Santos et al.~\cite{santos2026conflictairan} is a learned gated model, scored by its gate between a parameter
family and a target. INT is the intent-conflict statistic of PACIFISTA~\cite{PACIFISTA}, the distance between the
outcome distributions of accepted and rejected units. QACM~\cite{QACM2024} fits the KPI predictor of Wadud et al.\ on
the acted-on cell and scores the predicted change, so it never declares a neighbour or far edge. On the confounded logs,
each tuned baseline also received a threshold calibrated on the placebo logs themselves, which allows at most one
placebo edge and is not normally available to an operator (Table~\ref{tab:maps_full}).

\subsection{Ground truth and estimation targets}
\label{app:ko_labels}\claim{EMP-GT-LABELS}
\label{app:estimands}

Table~\ref{tab:system_hypotheses} gives the knockout label of each of the 60 hypotheses on the randomised logs, under
the rule of Section~\ref{sec:gt}. Every target is a 60~s mode contrast oriented by the request's direction, over a
90~s window (Section~\ref{sec:estimand}); Table~\ref{tab:estimands} lists the target of each kind of log and how it is
estimated.

\begin{table}[tbp]
\caption{Knockout ground truth, randomised logs.}
\label{tab:system_hypotheses}
% Tabular copied from tables/tab_system_hypotheses.tex (cddfig tab_system_hypotheses.py); parameter names glossed in
% the first column instead of the caption (Step 5 appendix rework).
\setlength{\tabcolsep}{3pt}
\begin{tabular}{@{}l l c c c c c@{}}
\toprule
\shortstack[l]{Parameter\\(writer)} & Relation & \shortstack{protected-slice\\violations} & \shortstack{service\\violations} & \shortstack{radio-link\\failures} & energy & \shortstack{served\\load} \\
\midrule
macro carrier & own & true $-$ & true $-$ & indeterminate & true $+$ & null \\
(ES) & neighbour & indeterminate & null & indeterminate & indeterminate & null \\
 & far & indeterminate & null & indeterminate & null & null \\
\midrule
pico sleep & own & indeterminate & indeterminate & true $-$ & null & true $-$ \\
(ES) & neighbour & true $+$ & true $+$ & true $+$ & true $+$ & true $+$ \\
 & far & true $+$ & true $+$ & true $+$ & true $+$ & true $+$ \\
\midrule
transmit power & own & true $+$ & true $+$ & true $+$ & true $+$ & true $+$ \\
(Power ES) & neighbour & true $-$ & true $-$ & true $-$ & true $-$ & true $-$ \\
 & far & indeterminate & true $-$ & indeterminate & true $-$ & true $-$ \\
\midrule
protected share & own & true $-$ & null & null & null & null \\
(Slice Guarantee) & neighbour & null & null & null & null & null \\
 & far & null & null & null & null & null \\
\bottomrule
\end{tabular}
\end{table}

\begin{table}[tbp]
\caption{Estimation targets.}
\label{tab:estimands}
\footnotesize
\setlength{\tabcolsep}{4pt}
\begin{tabular*}{\textwidth}{@{\extracolsep{\fill}}>{\raggedright}p{0.2\textwidth}>{\raggedright}p{0.27\textwidth}>{\raggedright}p{0.21\textwidth}>{\raggedright\arraybackslash}p{0.24\textwidth}@{}}
\toprule
Target & Other units in the window & Unit weighting & Estimated by \\
\midrule
$\tau_{\pi}$, randomised logs & random modes, $\pi_u = 0.5$ & equal & design slope (Section~\ref{sec:slope}) \\
$\tau_{\pi}$, confounded logs & random modes, incumbent propensities & $\pi_u(1 - \pi_u)$ & design slope \\
$\tau_{\mathrm{AA}}$, ground truth & all accepted & equal over sampled units & knockout rollouts \\
\botrule
\end{tabular*}
\end{table}

\subsection{Decision arms and metrics}
\label{app:arms}

Table~\ref{tab:arms} lists the arms of the two decision datasets and the decision metrics with their limits
(Section~\ref{sec:arms}). The recovery $R$ is the primary metric; an arm is eligible when its energy retention and all
four guard ratios are within their limits, and an ineligible arm is scored by $R^{*} = \min(R, 0)$.

\begin{table}[tbp]
\caption{Decision arms and metrics.}
\label{tab:arms}
\footnotesize
\setlength{\tabcolsep}{3pt}
\begin{tabular*}{\textwidth}{@{\extracolsep{\fill}}>{\raggedright}p{0.24\textwidth}>{\raggedright\arraybackslash}p{0.7\textwidth}@{}}
\toprule
Arm & Rule \\
\midrule
\multicolumn{2}{@{}l}{\textit{Static rules}} \\
Accept-all & Accepts every request; $R = 0$ by definition \\
Freeze & Rejects every request; energy retention 0 \\
Single-xApp subsets & Only the listed xApps run; the best one sets $R = 1$ \\
Never-sleep & Rejects pico-sleep requests; the best static rule on development seeds \\
Incumbent & The logging policy of the confounded logs \\
Static envelope & Rejects power-up, sleep and macro carrier-off after warm-up (without the gate only) \\
\midrule
\multicolumn{2}{@{}l}{\textit{Map-driven referees}} \\
PMRT map & PMRT's declared edges, with design-slope magnitudes \\
Oracle map & The true edges of the ground truth, with its means \\
Random map & PMRT's number of edges, at random keys and signs, with ground-truth magnitudes \\
No-map control & Assumes that energy-saving requests save own-cell energy and raise own and neighbour violations \\
Associational maps & Six baselines at two thresholds each, and Granger~+~BY (13 maps) \\
Ungated PMRT & The referee without the gate, rerun on the seeds of the gated dataset \\
Random rejection & Rejects each request with probability 0.25, 0.35 or 0.50 \\
\bottomrule
\end{tabular*}

\par\medskip

\begin{tabular*}{\textwidth}{@{\extracolsep{\fill}}>{\raggedright}p{0.3\textwidth}>{\raggedright}p{0.48\textwidth}>{\raggedright\arraybackslash}p{0.16\textwidth}@{}}
\toprule
Metric & Definition & Limit \\
\midrule
Protected violations $V$ & Per protected-UE hour, pooled over the seeds & Objective \\
Recovery $R$ & $(V_{\mathrm{AA}} - V) / (V_{\mathrm{AA}} - V_{\mathrm{ref}})$, with $V_{\mathrm{ref}}$ the best single-xApp subset & -- \\
Energy retention & Share of the energy-saving subset's saving over freeze & $\ge 0.90$ \\
Guard ratios & Service violations, non-protected eMBB and low-latency violations, RLFs; each over accept-all & $\le 1.10$ \\
Margin over associational maps & $R^{*}$ minus that of the best associational map & $\ge 0.10$; bound $> 0$ \\
Margin over never-sleep & $R^{*}$ minus that of never-sleep (gated dataset) & $> 0$; bound $> 0$ \\
RLF bound & Upper 90~\% bound of the RLF ratio (gated dataset) & $\le 1.20$ \\
\bottomrule
\end{tabular*}
\end{table}

\subsection{Controlled environments}
\label{app:envs}

Table~\ref{tab:study4_design} lists the five environments and the four ways in which their actions are logged
(Section~\ref{sec:envs}). A condition is one environment, logging design, confounding strength $\lambda$ and sample
size, with $n \in \{500, 1\,000, 4\,000, 8\,000, 24\,000\}$. A truth-null edge joins a real action to a KPI that it leaves
unaffected in the known graph; a placebo edge joins a KPI to an extra action column that is logged with a known design
but acts on nothing; and a confounded-null edge, in the latent-confounding environment only, joins a KPI to such a
column drawn by the confounded policy. Under unlogged confounding, the negative control, the real actions have no logged
design, so the design-based tests are read on the placebo edge only (Section~\ref{sec:study4_results}). The two PMRT
variants of these environments omit the Huber clip of the unit weight, which raised the null rate on development data;
the PMRT used in the simulator keeps the clip, and its validity rests on the checks of Section~\ref{sec:res_validity}.

\begin{table}[tbp]
\caption{Controlled environments and logging designs.}
\label{tab:study4_design}
% Tabular from tables/tab_study4_design.tex (cddfig tab_study4_design.py), with plain names for E1-E5 and R1-R4
% (Step 5 appendix rework; names as in fig_envs and fig_logging_validity).
\footnotesize
\setlength{\tabcolsep}{4pt}
\begin{tabular*}{\textwidth}{@{\extracolsep{\fill}}l p{5.6cm} c c@{}}
\toprule
Environment & Mechanism & True edges & Null edges \tabularnewline
\midrule
Linear graph & \raggedright Clean, identifiable linear graph & 4 & 16 \tabularnewline
Rare regime & \raggedright Edges active only in a rare regime & 16 & 38 \tabularnewline
Delayed harm & \raggedright Harm through lagged KPI cascades (depth 3) & 4 & 21 \tabularnewline
Latent confounding & \raggedright Associational sign wrong at large confounding strength; $\lambda \in \{0, 0.5, 1, 1.5\}$ & 1 & 1 \tabularnewline
Composed & \raggedright The three preceding mechanisms combined & 6 & 14 \tabularnewline
\bottomrule
\end{tabular*}

\par\medskip

\begin{tabular*}{\textwidth}{@{\extracolsep{\fill}}l p{8.6cm}@{}}
\toprule
Logging design & How actions are logged \tabularnewline
\midrule
Independent & \raggedright Every action redrawn independently per observation (reference) \tabularnewline
Dithered setpoint & \raggedright 20 setpoints within $\pm 25$~\% of the range around its midpoint, each held for $n/20$ observations, plus independent dither of $\pm 10$~\% of the range \tabularnewline
Confounded policy & \raggedright Latent confounding only: a behaviour policy on an observed confounder, with its probabilities logged \tabularnewline
Unlogged confounding & \raggedright Latent confounding only: the confounder is latent and nothing is logged (negative control) \tabularnewline
\bottomrule
\end{tabular*}
\end{table}
